%!TEX root = main.tex
% ============================================================
% 【格式参考】论文 LaTeX 骨架 —— 默认对齐国赛官方模板层级（写作前先定骨架）
% 写作思路 / 每部分怎么写、受阻怎么办 → 见 templates/paper_writing_guide.md（主件）
% 格式与版式规则（写作前必读：三线表/题注/小标题/防溢出等）→ 见 templates/paper_format_rules.md
% 优秀结构 / 高分要点 / 常见扣分点 → 见 references/paper_structure_reference.md
% 格式预设 / 排版自检（Overfull 修复、长公式技巧）→ 见 references/format_presets.md
% 用法：复制为 main.tex，先按官方论文模板文件 / 赛题通知核对下面层级，再按注填空。
%      写作请【详实、讲清论证链条、图表自解释】；复杂图（流程图等）请用户协助制作。
% 铁律：写稿即遵守版式规则（templates/paper_format_rules.md），别把排版问题留到收尾。
%      编译后扫 main.log 的 "Overfull \hbox" 并修复；收尾前做渲染抽查（见格式规则 §8）。
% ============================================================
\documentclass[12pt,a4paper]{article}

% ---- 基础包 ----
\usepackage[margin=2.5cm]{geometry}% 以官方模板边距为准（如国赛上下≥2.5cm）
\usepackage{amsmath,amssymb,amsfonts}
\usepackage{mathtools}% 增强公式（拆分/标注）
\usepackage{graphicx}
\usepackage{booktabs}% 三线表：\toprule/\midrule/\bottomrule，无竖线
\usepackage{tabularx}
\usepackage{array}
\usepackage{multirow}
\usepackage{siunitx}
\usepackage{enumitem}
\usepackage{caption}% 题注补充行（小字）用
\usepackage[hidelinks]{hyperref}% 无彩色链接框（勿改回 colorlinks）
\usepackage{indentfirst}

% ---- 中文排版（国赛中文用 xelatex；字体以官方模板为准）----
% \usepackage{xeCJK}
% \setCJKmainfont{SimSun}           % 正文宋体
% \setCJKsansfont{SimHei}           % 标题黑体
% 字号宏：\zihao{2}（二号）\zihao{5}（五号）等需配合 xeCJK/ctex 使用

% ---- 长公式/超宽工具（用于修复 Overfull \hbox）----
% 长公式/多 cases 禁止排一行：用 \begin{align} ... &= ... \\ ... &= ... \end{align} 分行堆叠；
% 或 \begin{equation}\begin{aligned} ... \end{aligned}\end{equation}；必要时 {\small ...}；
% 极端缩字 \resizebox{\textwidth}{!}{$...$}（慎用）。
% 行内禁止"无断点超长串"（区间 [a,b]/[c,d]、URL、长小数串）→ 改表格或分号分行表达。

% ---- 引用：二选一，取消注释其一 ----
% 方案 A：biblatex（推荐，处理中文/多源更稳）
% \usepackage[style=gb7714-2015,backend=biber]{biblatex}
% \addbibresource{refs.bib}
% 方案 B：natbib + bibtex
% \usepackage[numbers,sort&compress]{natbib}
% \bibliographystyle{plainnat}

\title{[论文题目]}
\author{[姓名1] [学号1]\\[姓名2] [学号2]\\[姓名3] [学号3]}
\date{\today}

\begin{document}

% ========== 摘要页（版式见格式规则 §2：题目居中大号粗体 / "摘要"居中 / 正文小字号） ==========
% 国赛惯例：电子版第一页 = 摘要页，无参赛者信息；标题不靠 \maketitle 的默认排版，
% 用居中控制手动排：\begin{center}\zihao{2} 题目\end{center}（中文配合 xeCJK/ctex 的字号命令）。
\maketitle
\begin{abstract}
[一句话点题：针对什么背景/问题。]
[本工作做了什么：建模思路、采用的关键模型/算法（2-3 步）。]
[主要结果：关键数值/结论，1-3 个亮点。]
[结论/评价：模型有效性、稳健性、适用性。]
\textbf{关键词：}[关键词1]；[关键词2]；[关键词3]；[关键词4]；[关键词5]
\end{abstract}
% 注：摘要不要放"本文亮点/创新点"等非标准内容；删之。

% ========== 1 问题重述（凝练，不是抄题；说清"要交付什么"） ==========
\section{问题重述}
[清晰重述题目要求与交付物。]

% ========== 2 问题分析（问题X分析，逐问讲思路） ==========
\section{问题分析}
\subsection{问题一分析}
[问题一属于什么模型类型、可用什么算法、方向、难点。]
\subsection{问题二分析}
[……]
\subsection{问题三分析}
[……]
\subsection{问题四分析}
[……]

% ========== 3 模型假设（明确、可检验） ==========
\section{模型假设}
\begin{enumerate}
    \item [假设1：……（不成立的后果：……）]
    \item [假设2：……]
\end{enumerate}

% ========== 4 符号说明（与模型假设分节；单对一行三线表，见格式规则 §4） ==========
\section{符号说明}
% 符号表：每行一对（符号|含义|单位），不一行两排；无竖线三线表。
\begin{table}[htbp]
\centering
\begin{tabular}{clc}
\toprule
符号 & 含义 & 单位 \\
\midrule
$n$ & [含义] & [单位] \\
$c$ & [含义] & [单位] \\
\bottomrule
\end{tabular}
\end{table}

% ========== 5 模型的建立与求解（问题X 嵌在此节下，勿把每问做成顶层节） ==========
\section{模型的建立与求解}

\subsection{问题一模型的建立与求解}
\subsubsection*{模型的建立}
[为什么这样建模（一段思路）→ 变量/参数/约束/目标；公式用 equation/align。]
\subsubsection*{模型的求解}
[求解方法、算法步骤；长公式/多 cases 用 align 分行堆叠。]
\subsubsection*{结果分析}
[结果展示：三线表 + 图；文字直接接图表陈述数据与结论（不另开"结果解读"）。]
% 细节说明用 \paragraph{…}（未编号、嵌套正确）或段落内粗体引导，不堆同级 \subsubsection。
% 图：\begin{figure}[htbp]\centering\includegraphics[width=0.6\textwidth]{fig.pdf}\caption{一行主标题}\caption*{补充说明：坐标轴/读法（小字，caption 包支持）}\end{figure}

\subsection{问题二模型的建立与求解}
[…… 同上结构：模型的建立 / 模型的求解 / 结果分析]

\subsection{问题三模型的建立与求解}
[……]

\subsection{问题四模型的建立与求解}
[……]

% ========== 6 模型的分析与检验（假设/误差/敏感性/稳健性汇总） ==========
\section{模型的分析与检验}
[误差来源、实际信度（如实际 α/β、解析 vs 蒙特卡洛偏差）、敏感性/稳健性结论。]

% ========== 7 模型的评价、改进与推广 ==========
\section{模型的评价、改进与推广}
\subsection{优点}
[模型/算法优点。]
\subsection{缺点与局限}
[局限、假设不成立时的风险（诚实）。]
\subsection{改进与推广}
[可扩展方向。]

% ========== 参考文献（\cite{} ↔ .bib 一一对应） ==========
% 方案 A：biblatex
% \printbibliography
% 方案 B：natbib + bibtex
% \bibliography{refs}

% ========== 附录 ==========
\appendix
\section{附录：主要代码 / 支撑材料}
[按赛题要求放核心代码、数据表、中间推导等；程序清单类表格用固定列宽+小字防超宽。]

\end{document}

% ============================================================
% 【美赛（US_MCM/ICM）变体 · 仅作注释，勿与上方同时启用】
% 结构：Summary(单页) → 1 Introduction → 2 Assumptions → 3 Notation
%       → 4 Model (每问) → 5 Results & Analysis → 6 Sensitivity/Robustness
%       → 7 Conclusion → References → Appendix；全英文；Team # 标注。
% ============================================================
